\chapter{Planificación y recursos}
\label{ch:planificacion_recursos}

\section{Enfoque de planificación}
\label{sec:enfoque_planificacion}

La planificación de Harmony Hub se ha organizado mediante una adaptación ligera
de Scrum al contexto de un TFG. No se ha aplicado Scrum como marco empresarial
formal, sino como metodología de trabajo incremental para ordenar el backlog,
priorizar entregables, revisar avances con frecuencia y reducir el riesgo de
construir funcionalidades sin validación progresiva.

Esta decisión resulta coherente con la naturaleza del proyecto por tres
motivos. En primer lugar, Harmony Hub combina frontend móvil, backend, modelo
supervisado, persistencia en la nube e integración con un servicio externo como
Spotify, por lo que era preferible avanzar por incrementos verificables en vez
de seguir una planificación rígida de una sola pasada. En segundo lugar, parte
del comportamiento final dependía de pruebas reales y de la evolución del flujo
de uso, especialmente en la recomendación y en la interfaz. En tercer lugar, el
TFG exigía compaginar desarrollo técnico, validación y documentación, de modo
que una organización iterativa facilitaba equilibrar implementación y memoria.

Como soporte práctico de esta organización se utilizó también un
\href{https://trello.com/invite/b/69f0d52f02062b866ef2c7b1/ATTI8242c2e81725ea5b57b648b8e983e91652F5E110/harmony-hub}{tablero de Trello del proyecto}.
Ese tablero actuó como vista operativa del backlog y del avance cotidiano,
mientras que la memoria recoge la versión consolidada y argumentada de esa
misma planificación.

\section{Adaptación de Scrum al TFG}
\label{sec:adaptacion_scrum_tfg}

La adaptación realizada conserva la lógica básica de Scrum, pero simplificada
para un proyecto individual. En este capítulo, los términos principales se usan
con el siguiente significado práctico: backlog de producto como lista
priorizada de tareas, sprint como iteración corta de trabajo y
sprint planning como selección del alcance inmediato de cada iteración.

\begin{itemize}
    \item \textbf{Backlog de producto}. Se definió una lista inicial priorizada con historias y tareas agrupadas en bloques: autenticación, integración con Spotify, check-in, motor de recomendación, feedback, aprendizaje, seguridad de datos, rediseño de interfaz y documentación.
    \item \textbf{Sprints cortos}. El trabajo se estructuró en iteraciones de aproximadamente dos semanas, suficientes para construir incrementos reales y revisarlos antes de pasar al siguiente bloque.
    \item \textbf{Planificación de sprint}. Al inicio de cada iteración se seleccionaban objetivos concretos y un alcance limitado, priorizando siempre el flujo principal antes que las extensiones.
    \item \textbf{Revisión continua}. Cada incremento se validó mediante compilación, ejecución manual de la app o comprobación de logs del backend.
    \item \textbf{Retrospectiva informal}. Tras cada bloque se registraron incidencias, limitaciones detectadas y cambios de prioridad, por ejemplo al reforzar el motor heurístico, revisar reglas de Firestore o rediseñar pantallas completas.
\end{itemize}

Al tratarse de un proyecto individual, las funciones de priorización,
seguimiento y desarrollo quedaron concentradas en una sola persona. Aun así, se
mantuvo el valor metodológico del enfoque iterativo: backlog priorizado,
iteraciones, validación continua y entregables incrementales.

En la práctica, el tablero de Trello permitió trasladar esta adaptación a un
formato visual sencillo, útil para reagrupar tareas, mover bloques entre
estados y mantener visible la relación entre backlog, incidencias y trabajo en
curso.

\section{Backlog funcional priorizado}
\label{sec:backlog_funcional_priorizado}

La Tabla~\ref{tab:backlog_priorizado} resume los grandes bloques del backlog
ordenados por prioridad funcional.

\begin{table}[H]
\centering
\scriptsize
\renewcommand{\arraystretch}{1.1}
\setlength{\tabcolsep}{3pt}
\begin{tabularx}{\textwidth}{|P{1.45cm}|P{3.55cm}|P{2.45cm}|Y|}
\hline
\textbf{Prioridad} & \textbf{Bloque de trabajo} & \textbf{Valor principal} & \textbf{Justificación} \\
\hline
Alta & Autenticación y acceso inicial & Entrada al sistema & Sin acceso autenticado no existe flujo funcional utilizable. \\
\hline
Alta & Integración con Spotify & Generación real de playlists & La propuesta de valor del proyecto depende de conectar cuenta, consultar perfil y crear playlists reales. \\
\hline
Alta & Check-in contextual & Captura de entrada & El motor de recomendación necesita inputs de estado emocional, objetivo y contexto para operar. \\
\hline
Alta & Motor de recomendación heurístico & Núcleo funcional mínimo & Permite generar sesiones útiles incluso antes de disponer de un modelo supervisado maduro. \\
\hline
Media & Persistencia, historial y feedback & Trazabilidad y aprendizaje & Cierran el ciclo de uso y habilitan la recopilación de ejemplos para aprendizaje posterior. \\
\hline
Media & Componente supervisado de sesión & Personalización incremental & Mejora la priorización de modos, pero no bloquea la operatividad del sistema. \\
\hline
Media & Seguridad y privacidad del dato & Viabilidad ética y técnica & El proyecto trata información emocional y requiere separación entre bloque público y cifrado. \\
\hline
Media & Rediseño integral de interfaz & Calidad de experiencia & Relevante para la presentación final y para reducir fricción en el uso. \\
\hline
Baja & Optimización fina y mejoras estéticas secundarias & Refinamiento & Se abordaron después de asegurar el flujo principal y la consistencia funcional. \\
\hline
\end{tabularx}
\caption{Backlog funcional priorizado durante el desarrollo de Harmony Hub.}
\label{tab:backlog_priorizado}
\end{table}

\section{Fases de desarrollo}
\label{sec:fases_desarrollo}

Aunque la ejecución práctica se organizó mediante iteraciones cortas, el
proyecto puede describirse también a través de una secuencia de fases que ayuda
a interpretar su evolución global. Esta lectura complementa la planificación
por sprints y permite relacionar con mayor claridad cada bloque de trabajo con
el tipo de resultado perseguido.

\begin{itemize}
    \item \textbf{Preparación técnica y estudio previo}. Se definió el alcance inicial, se revisó la viabilidad tecnológica y se estableció la base de trabajo con Flutter, FastAPI, Firebase y Spotify.
    \item \textbf{Análisis y especificación}. Se concretaron requisitos, actores, casos de uso, necesidades de persistencia y criterios de privacidad para el tratamiento del dato emocional.
    \item \textbf{Diseño del sistema}. Se modelaron la arquitectura general, la separación entre cliente y backend, las colecciones documentales de Firestore, los flujos principales y la lógica híbrida de recomendación.
    \item \textbf{Implementación funcional}. Se desarrollaron el acceso autenticado, la conexión con Spotify, el check-in contextual, la recomendación de sesión, la generación de playlists, el historial y el cierre con feedback.
    \item \textbf{Aprendizaje y validación}. Se incorporaron la capa de aprendizaje individual, el conjunto supervisado de sesión, el mantenimiento automático del modelo y las pruebas funcionales, técnicas y longitudinales.
    \item \textbf{Refinamiento y documentación}. Finalmente se ajustaron interfaz, trazabilidad, consistencia entre módulos y memoria, con el fin de consolidar una versión entregable y coherente del sistema.
\end{itemize}

\section{Planificación por iteraciones}
\label{sec:planificacion_iteraciones}

La ejecución del proyecto se organizó en una secuencia de iteraciones con un
objetivo principal por sprint. La Tabla~\ref{tab:sprints_harmony_hub} recoge la
estructura general seguida.

\begin{table}[H]
\centering
\small
\renewcommand{\arraystretch}{1.1}
\begin{tabularx}{\textwidth}{|P{1.7cm}|P{2.4cm}|Y|P{2.7cm}|}
\hline
\textbf{Sprint} & \textbf{Meta principal} & \textbf{Entregables o tareas destacadas} & \textbf{Resultado esperado} \\
\hline
Sprint 0 & Definición y base técnica & Análisis del problema, objetivos, selección tecnológica, preparación de repositorio, estructura Flutter/FastAPI y configuración inicial de Firebase. & Base de trabajo estable y backlog inicial definido. \\
\hline
Sprint 1 & Flujo de acceso y conexión & Registro, inicio de sesión, recuperación de contraseña, conexión con Spotify y validación de cuenta enlazada. & Entrada real al sistema y acceso a datos del perfil musical. \\
\hline
Sprint 2 & Check-in y flujo principal & Captura de inputs, medición opcional del entorno, persistencia inicial y navegación hacia recomendación. & Flujo contextual desde la entrada emocional hasta la propuesta de sesión. \\
\hline
Sprint 3 & Motor de recomendación e integración backend & Perfiles de generación, búsquedas, ranking heurístico, creación de playlist y persistencia de sesión. & Generación funcional de playlists reales en Spotify. \\
\hline
Sprint 4 & Feedback, historial y aprendizaje & Cierre de sesión con feedback, historial, modos guardados, aprendizaje del usuario y ejemplos para entrenamiento. & Ciclo completo de uso y base de datos de sesiones consolidada. \\
\hline
Sprint 5 & Modelo supervisado y validación & Entrenamiento del clasificador de sesión, calidad del modelo, revisión de reglas, pruebas funcionales y correcciones. & Personalización progresiva validada y estado técnico verificado. \\
\hline
Sprint 6 & Rediseño y cierre documental & Rediseño de pantallas clave, coherencia visual, ajuste de flujos, memoria y revisión final. & Prototipo final presentable y documentación alineada con la implementación real. \\
\hline
\end{tabularx}
\caption{Planificación del proyecto por iteraciones inspiradas en Scrum.}
\label{tab:sprints_harmony_hub}
\end{table}

\section{Cronograma resumido}
\label{sec:cronograma}

Para esta memoria se documenta el cronograma en formato tabular, ya que los
diagramas temporales se reservarán para una fase posterior en la que el texto y
los bloques del proyecto queden completamente estabilizados. La
Tabla~\ref{tab:cronograma_resumido_scrum} resume la distribución temporal
orientativa de las iteraciones.

\begin{table}[H]
\centering
\small
\renewcommand{\arraystretch}{1.1}
\begin{tabularx}{\textwidth}{|P{1.9cm}|P{2.4cm}|Y|}
\hline
\textbf{Iteración} & \textbf{Duración orientativa} & \textbf{Bloque de trabajo dominante} \\
\hline
Sprint 0 & 1 semana & Definición del alcance, arquitectura inicial y preparación del entorno. \\
\hline
Sprint 1 & 2 semanas & Acceso, autenticación y conexión con Spotify. \\
\hline
Sprint 2 & 2 semanas & Check-in contextual, navegación y persistencia inicial. \\
\hline
Sprint 3 & 2 semanas & Motor heurístico, backend y creación de playlists. \\
\hline
Sprint 4 & 2 semanas & Feedback, historial, aprendizaje y datos de entrenamiento. \\
\hline
Sprint 5 & 2 semanas & Entrenamiento del modelo, pruebas y correcciones. \\
\hline
Sprint 6 & 2 semanas & Rediseño de interfaz, revisión final y documentación. \\
\hline
\end{tabularx}
\caption{Cronograma resumido del proyecto siguiendo una estructura iterativa.}
\label{tab:cronograma_resumido_scrum}
\end{table}

Más que fijar fechas cerradas, esta planificación ha servido para mantener una
secuencia lógica: primero asegurar la viabilidad funcional mínima, después
añadir trazabilidad y aprendizaje, y por último refinar calidad visual,
robustez y documentación.

\section{Desviaciones respecto a la planificación inicial}
\label{sec:desviaciones_planificacion}

La planificación inicial proporcionó una guía útil, pero la evolución real del
proyecto exigió introducir varios ajustes para mantener la coherencia entre
ambición funcional, viabilidad técnica y calidad final del entregable. Estas
desviaciones no alteraron el propósito general de Harmony Hub, aunque sí
reordenaron prioridades y forzaron decisiones de diseño más prudentes en
algunos bloques.

Una primera desviación relevante afectó al papel de Spotify dentro del sistema.
En un planteamiento temprano se contemplaba una dependencia mayor de búsquedas y
afinidad dinámica sobre el catálogo externo. La experiencia de integración, las
limitaciones de materialización y la necesidad de preservar la trazabilidad
llevaron a reforzar un catálogo propio como base principal del ranking musical,
reservando Spotify para autenticación, contexto musical básico y creación final
de playlists reproducibles.

También se produjeron ajustes en persistencia e historial. La necesidad de
separar bloque público y bloque privado, alinear reglas de Firestore con los
documentos generados y mantener un historial consistente con los datos
persistidos obligó a refinar la estructura documental y algunos endpoints del
backend. Este trabajo añadió complejidad respecto a la previsión inicial, pero
mejoró la coherencia técnica y la protección del dato personal.

En la capa de aprendizaje, la desviación principal consistió en abandonar una
lectura simplificada del progreso basada solo en acumulación de sesiones. La
evidencia longitudinal mostró que era más riguroso distinguir entre aprendizaje
individual del perfil y activación del clasificador supervisado global. Como
consecuencia, la solución adoptó una política de puertas de calidad y de
abstención explícita, lo que reorientó parte del trabajo de validación hacia la
interpretación del balance entre clases y la madurez específica por mood.

Por último, el rediseño de interfaz y la revisión documental adquirieron más
peso del previsto al final del proyecto. Lejos de constituir un añadido
superficial, ambos bloques resultaron necesarios para que la solución
expresara con claridad la lógica de aprendizaje, la trazabilidad de uso y la
separación entre lo ya maduro y lo que aún permanece en validación prudente.

\section{Seguimiento y control del avance}
\label{sec:seguimiento_control_avance}

El seguimiento del avance se ha apoyado en cuatro mecanismos sencillos pero
efectivos:

\begin{itemize}
    \item \textbf{Criterios de terminado}. Cada sprint se consideraba cerrado cuando el incremento compilaba, podía recorrerse funcionalmente y quedaba documentado a nivel de código o memoria.
    \item \textbf{Validación técnica continua}. Se utilizaron comprobaciones de análisis estático del cliente, compilación del backend, revisión de logs y compilación de la memoria.
    \item \textbf{Repriorización}. Cuando aparecían limitaciones relevantes, como restricciones de Spotify o problemas de permisos en Firestore, el backlog se reajustaba para corregir primero el bloqueo.
    \item \textbf{Integración con la memoria}. La documentación no se dejó para el final, sino que se fue actualizando conforme se estabilizaban arquitectura, pruebas y decisiones de diseño.
\end{itemize}

Este enfoque ha permitido que la planificación no quedara separada de la
realidad técnica del proyecto. En lugar de medir solo horas o tareas, se han
cerrado incrementos funcionales verificables.

Además, el tablero de Trello se utilizó como mecanismo auxiliar de control del
avance. Su función principal no fue sustituir la planificación formal, sino
servir como soporte visual para revisar prioridades, detectar bloqueos y
mantener coherencia entre el trabajo diario y la estructura por iteraciones
documentada en este capítulo.

\section{Recursos empleados}
\label{sec:recursos_empleados}

Los recursos empleados pueden agruparse en cuatro categorías: recursos
humanos, recursos hardware, recursos software y servicios externos. Esta
separación permite describir mejor qué papel ha desempeñado cada elemento
dentro de la ejecución real del proyecto.

\subsection*{Recursos humanos}

\begin{table}[H]
\centering
\small
\renewcommand{\arraystretch}{1.1}
\begin{tabularx}{\textwidth}{|P{3.2cm}|Y|}
\hline
\textbf{Elemento} & \textbf{Uso dentro del proyecto} \\
\hline
Responsable del desarrollo & Análisis, diseño, implementación, pruebas, validación y documentación integral del sistema. \\
\hline
Tutor académico & Supervisión metodológica, revisión del enfoque y orientación general del trabajo. \\
\hline
\end{tabularx}
\caption{Recursos humanos empleados en el proyecto.}
\label{tab:recursos_humanos}
\end{table}

\subsection*{Recursos hardware}

\begin{table}[H]
\centering
\small
\renewcommand{\arraystretch}{1.1}
\begin{tabularx}{\textwidth}{|P{3.2cm}|Y|}
\hline
\textbf{Elemento} & \textbf{Uso dentro del proyecto} \\
\hline
Ordenador personal de desarrollo & Programación, entrenamiento del modelo, compilación de la memoria y ejecución local del backend. \\
\hline
Dispositivo móvil y emulador Android & Pruebas funcionales de la app Flutter, validación visual de la interfaz y comprobación del flujo real de uso. \\
\hline
\end{tabularx}
\caption{Recursos hardware empleados en el proyecto.}
\label{tab:recursos_hardware}
\end{table}

\subsection*{Recursos software}

\begin{table}[H]
\centering
\small
\renewcommand{\arraystretch}{1.1}
\begin{tabularx}{\textwidth}{|P{3.2cm}|Y|}
\hline
\textbf{Elemento} & \textbf{Uso dentro del proyecto} \\
\hline
Flutter y Dart & Desarrollo del cliente móvil multiplataforma. \\
\hline
FastAPI y Python & Implementación del backend, servicios y scripts de entrenamiento. \\
\hline
Firebase / Firestore & Autenticación, persistencia de datos, reglas de acceso y soporte documental del historial. \\
\hline
scikit-learn & Entrenamiento, validación y auditoría del modelo supervisado de sesión. \\
\hline
Trello & Seguimiento visual del backlog, de las iteraciones y del avance diario. \\
\hline
\LaTeX & Redacción y maquetación de la memoria final. \\
\hline
\end{tabularx}
\caption{Recursos software empleados en el proyecto.}
\label{tab:recursos_software}
\end{table}

\subsection*{Servicios externos}

\begin{table}[H]
\centering
\small
\renewcommand{\arraystretch}{1.1}
\begin{tabularx}{\textwidth}{|P{3.2cm}|Y|}
\hline
\textbf{Elemento} & \textbf{Uso dentro del proyecto} \\
\hline
Spotify Web API & Autenticación musical, acceso al perfil básico, búsqueda de coincidencias y creación de playlists reproducibles. \\
\hline
Cuenta de Spotify de prueba & Validación del flujo real de generación de playlists. \\
\hline
Proyecto Firebase del TFG & Persistencia documental, reglas de seguridad y trazabilidad del uso. \\
\hline
\end{tabularx}
\caption{Servicios externos empleados en el proyecto.}
\label{tab:recursos_servicios}
\end{table}
